% ============================================================
% Week 1 — 02 Mar – 06 Mar 2026
% Compile:  cd report && make week_01
% ============================================================
\documentclass[11pt, a4paper]{article}

\usepackage[margin=2.5cm]{geometry}
\usepackage{graphicx}
\usepackage{booktabs}
\usepackage{multirow}
\usepackage{array}
\usepackage{xcolor}
\usepackage{hyperref}
\usepackage{amsmath}
\usepackage{caption}
\usepackage{subcaption}
\usepackage{float}
\usepackage{enumitem}
\usepackage{fancyhdr}
\usepackage{titlesec}
\usepackage{parskip}

\definecolor{sota}{HTML}{2E7D32}
\definecolor{good}{HTML}{1565C0}
\definecolor{warn}{HTML}{E65100}

\setlength{\headheight}{14pt}
\pagestyle{fancy}
\fancyhf{}
\lhead{\textbf{ESS Research} --- Week 1 Update}
\rhead{Annotation-Free Surgical Instrument Segmentation}
\cfoot{\thepage}

\titleformat{\section}{\large\bfseries}{}{0em}{}[\titlerule]
\titleformat{\subsection}{\normalsize\bfseries}{}{0em}{}

\hypersetup{colorlinks=true, linkcolor=blue, urlcolor=blue}
\newcolumntype{C}[1]{>{\centering\arraybackslash}p{#1}}

\begin{document}

% ── Cover ────────────────────────────────────────────────────
\begin{center}
    {\LARGE \textbf{Towards Annotation-Free Surgical Instrument Segmentation}}\\[0.3em]
    {\large via Foundation Model Adaptation and Robot Kinematics}\\[0.5em]
    \rule{0.5\linewidth}{0.4pt}\\[0.4em]
    {\normalsize Weekly Research Progress Report --- \textbf{Week 1}}\\[0.3em]
    {\normalsize Week dates: 02 Mar -- 06 Mar 2026 \quad|\quad Report date: 06 March 2026}\\[0.3em]
    {\normalsize \href{https://github.com/skhan61/ess-research}{\texttt{github.com/skhan61/ess-research}}}
\end{center}
\vspace{0.5em}

\noindent This work investigates a three-phase path toward a \textbf{deployable,
annotation-free surgical instrument segmentation system} using SAM3 as the foundation
model backbone. \textbf{Phase~1}: SAM3+LoRA vs.\ task-specific SOTA using GT-derived
prompts. \textbf{Phase~2}: replace GT prompts with robot-kinematic-derived spatial
priors. \textbf{Phase~3}: feed SAM3 output back into the robot control loop.
\textit{This report covers Phase~1 (Week~1).}

\hrule
\vspace{1em}

% ============================================================
\section{Project Overview}
% ============================================================

\subsection{Objective}
Investigate whether \textbf{SAM3} (\texttt{facebook/sam3}) can match or exceed
HFRF-Net on surgical instrument segmentation using lightweight \textbf{LoRA
fine-tuning} on the ViT image encoder attention projections.

\subsection{Dataset}
The \textbf{UW-Sinus-Surgery-CL} dataset contains paired endoscopic video frames
and binary instrument masks from two domains:

\begin{center}
\begin{tabular}{llll}
\toprule
\textbf{Domain} & \textbf{Sessions} & \textbf{Frames} & \textbf{Characteristics} \\
\midrule
Cadaver & S01--S10 & 4{,}345 & Clean anatomy, consistent lighting \\
Live    & L01--L03 & 4{,}658 & Blood, smoke, motion blur \\
\midrule
\textbf{Total} & & \textbf{9{,}003} & \\
\bottomrule
\end{tabular}
\end{center}

\subsection{Model Architecture}
SAM3 has three components:
\begin{enumerate}[noitemsep]
    \item \textbf{Image Encoder} --- ViT, patch size 14; input $336\times336$.
          LoRA on \texttt{q\_proj}/\texttt{v\_proj}, rank $r=4$, $\alpha=16$,
          dropout $0.1$. Only $\approx0.5\%$ of weights trainable.
    \item \textbf{Prompt Encoder} --- encodes text, box, and point prompts.
          SAM3 always requires \texttt{input\_ids} (text).
    \item \textbf{Mask Decoder} --- cross-attends image + prompt embeddings,
          outputs logits $(B,1,H,W)$.
\end{enumerate}

\subsection{Training Configuration}
\begin{center}
\begin{tabular}{ll}
\toprule
\textbf{Parameter} & \textbf{Value} \\
\midrule
Optimizer        & AdamW, lr $1\times10^{-4}$ \\
Loss             & Dice + BCE (MONAI \texttt{DiceCELoss}) \\
Max epochs       & 50, early stopping patience = 10 \\
Batch size       & 8, image size $336\times336$ \\
Precision        & \texttt{16-mixed} (LoRA), \texttt{32-true} (zero-shot) \\
\bottomrule
\end{tabular}
\end{center}

% ============================================================
\section{Experimental Plan}
% ============================================================

Entry point: \texttt{uv run python scripts/run\_all\_experiments.py --plan research\_plan.yaml}

\begin{center}
\begin{tabular}{C{1cm} C{3.2cm} C{5cm} C{4cm}}
\toprule
\textbf{ID} & \textbf{Title} & \textbf{Hypothesis} & \textbf{Experiments} \\
\midrule
H1 & Zero-shot Cadaver    & SAM3 over-segments; establishes baseline & C$\to$C (1 run) \\
H2 & Zero-shot Live (CV)  & Live harder; high variance               & L$\to$L (3-fold CV) \\
H3 & LoRA Within-domain   & LoRA corrects over-segmentation          & C$\to$C, L$\to$L $\times$3 \\
H4 & Prompt Ablation      & box $>$ all $\approx$ point $>$ text     & box/point/all on C$\to$C + L$\to$L $\times$3 \\
H5 & Consolidated Re-eval & Checkpoint determinism verification       & All H3+H4 re-tested \\
\bottomrule
\end{tabular}
\end{center}

% ============================================================
\section{Results}
% ============================================================

\subsection{H1 --- Zero-shot: Cadaver (C$\to$C)}

\begin{center}
\begin{tabular}{ccccc}
\toprule
\textbf{Experiment} & \textbf{Dice} & \textbf{IoU} & \textbf{Precision} & \textbf{Recall} \\
\midrule
C$\to$C fold1 & 0.6016 & 0.4302 & \textcolor{warn}{0.4447} & \textcolor{sota}{0.9296} \\
\bottomrule
\end{tabular}
\end{center}

\textbf{Observation:} Classic over-segmentation --- recall 0.93 (finds instruments)
but precision 0.44 (blobs extend beyond boundaries). Dice $=0.60$ confirms the model
is sensing instruments but not their exact shape. Motivates LoRA fine-tuning.

\begin{figure}[H]
    \centering
    \includegraphics[width=\linewidth,height=10cm,keepaspectratio]{../../outputs/H1/C_C_fold1/predictions/predictions.eps}
    \caption{H1: Zero-shot cadaver. Each row: input, GT mask, predicted mask.
             Over-segmentation is clearly visible.}
\end{figure}

\subsection{H2 --- Zero-shot: Live Surgery (L$\to$L, 3-fold CV)}

\begin{center}
\begin{tabular}{ccccc}
\toprule
\textbf{Fold} & \textbf{Dice} & \textbf{IoU} & \textbf{Precision} & \textbf{Recall} \\
\midrule
L$\to$L fold1 & 0.5679 & 0.3966 & 0.4279 & 0.8442 \\
L$\to$L fold2 & \textcolor{warn}{0.4422} & 0.2839 & \textcolor{warn}{0.2966} & 0.8691 \\
L$\to$L fold3 & \textcolor{good}{0.7185} & 0.5606 & 0.5934 & 0.9104 \\
\midrule
\textbf{Mean $\pm$ std} & $0.576 \pm 0.138$ & & & \\
\bottomrule
\end{tabular}
\end{center}

\textbf{Observation:} Live is harder on average ($0.576 < 0.602$). Fold2 ($0.44$) is
the hardest session. Fold3 ($0.72$) is anomalously easy --- likely cleaner frames.
High variance ($\pm0.138$) shows session-specific difficulty dominates over any
systematic model limitation.

\begin{figure}[H]
    \centering
    \begin{subfigure}[b]{0.32\linewidth}
        \includegraphics[width=\linewidth]{../../outputs/H2/L_L_fold1/predictions/predictions.eps}
        \caption{Fold 1 --- Dice 0.568}
    \end{subfigure}
    \begin{subfigure}[b]{0.32\linewidth}
        \includegraphics[width=\linewidth]{../../outputs/H2/L_L_fold2/predictions/predictions.eps}
        \caption{Fold 2 --- Dice 0.442}
    \end{subfigure}
    \begin{subfigure}[b]{0.32\linewidth}
        \includegraphics[width=\linewidth]{../../outputs/H2/L_L_fold3/predictions/predictions.eps}
        \caption{Fold 3 --- Dice 0.719}
    \end{subfigure}
    \caption{H2: Zero-shot live surgery across three folds.}
\end{figure}

\subsection{H3 --- LoRA Fine-tuning: Within-domain}

\begin{center}
\begin{tabular}{ccccc}
\toprule
\textbf{Experiment} & \textbf{Dice} & \textbf{IoU} & \textbf{Precision} & \textbf{Recall} \\
\midrule
C$\to$C fold1   & \textcolor{sota}{\textbf{0.9329}} & 0.8742 & 0.9358 & 0.9300 \\
L$\to$L fold1   & 0.9049 & 0.8263 & 0.9033 & 0.9065 \\
L$\to$L fold2   & 0.8389 & 0.7225 & 0.8297 & 0.8482 \\
L$\to$L fold3   & 0.8886 & 0.7996 & 0.8623 & 0.9166 \\
\midrule
\textbf{L$\to$L mean} & $0.877 \pm 0.034$ & & & \\
\midrule
HFRF-Net (UW-Sinus live, SOTA) & 0.9374 & --- & --- & --- \\
\bottomrule
\end{tabular}
\end{center}

\textbf{Observation:} LoRA completely fixes over-segmentation. Precision jumps from
$0.44\to0.94$ on cadaver. C$\to$C Dice $0.9329$ is within \textbf{0.005 of SOTA}
(HFRF-Net was evaluated on the live subset; cadaver results were not reported
separately). Precision and recall are now balanced (within $0.02$), confirming
LoRA calibrated boundaries without sacrificing recall.

\begin{figure}[H]
    \centering
    \begin{subfigure}[b]{0.48\linewidth}
        \includegraphics[width=\linewidth]{../../outputs/H3/C_C_fold1/predictions/predictions.eps}
        \caption{C$\to$C fold1 --- Dice \textbf{0.9329}}
    \end{subfigure}
    \begin{subfigure}[b]{0.48\linewidth}
        \includegraphics[width=\linewidth]{../../outputs/H3/L_L_fold1/predictions/predictions.eps}
        \caption{L$\to$L fold1 --- Dice 0.9049}
    \end{subfigure}
    \caption{H3: LoRA fine-tuned predictions. Boundaries are tight vs.\ H1/H2.}
\end{figure}

\subsection{H4 --- Prompt Type Ablation}

\subsubsection{C$\to$C}
\begin{center}
\begin{tabular}{lcccc}
\toprule
\textbf{Prompt} & \textbf{Dice} & \textbf{IoU} & \textbf{Precision} & \textbf{Recall} \\
\midrule
text (H3 baseline) & 0.9329 & 0.8742 & 0.9358 & 0.9300 \\
point              & 0.9281 & 0.8658 & 0.9206 & 0.9357 \\
\textbf{box}       & \textcolor{sota}{\textbf{0.9476}} & \textbf{0.9004} & \textbf{0.9499} & \textbf{0.9453} \\
all                & 0.9453 & 0.8963 & 0.9490 & 0.9417 \\
\bottomrule
\end{tabular}
\end{center}

\subsubsection{L$\to$L (all three folds)}
\begin{center}
\begin{tabular}{llcccc}
\toprule
\textbf{Fold} & \textbf{Prompt} & \textbf{Dice} & \textbf{IoU} & \textbf{Prec.} & \textbf{Rec.} \\
\midrule
\multirow{4}{*}{fold1}
  & text  & 0.9049 & 0.8263 & 0.9033 & 0.9065 \\
  & point & 0.9132 & 0.8403 & 0.9075 & 0.9190 \\
  & \textbf{box} & \textcolor{sota}{\textbf{0.9593}} & \textbf{0.9218} & \textbf{0.9586} & \textbf{0.9601} \\
  & all   & 0.9580 & 0.9193 & 0.9630 & 0.9530 \\
\midrule
\multirow{4}{*}{fold2}
  & text  & 0.8389 & 0.7225 & 0.8297 & 0.8482 \\
  & point & 0.8766 & 0.7804 & 0.8447 & 0.9110 \\
  & \textbf{box} & \textcolor{sota}{\textbf{0.9320}} & \textbf{0.8726} & \textbf{0.9181} & \textbf{0.9462} \\
  & all   & 0.9327 & 0.8738 & 0.9174 & 0.9485 \\
\midrule
\multirow{4}{*}{fold3}
  & text  & 0.8886 & 0.7996 & 0.8623 & 0.9166 \\
  & point & 0.9007 & 0.8193 & 0.8832 & 0.9189 \\
  & \textbf{box} & \textcolor{sota}{\textbf{0.9548}} & \textbf{0.9134} & \textbf{0.9497} & \textbf{0.9598} \\
  & all   & 0.9523 & 0.9089 & 0.9478 & 0.9569 \\
\bottomrule
\end{tabular}
\end{center}

\textbf{Key observations:}
\begin{itemize}[noitemsep]
    \item \textbf{Box is the strongest prompt} across all experiments.
    \item \textbf{Box on L$\to$L surpasses HFRF-Net ($0.9374$)} in all three folds
          ($0.9593$, $0.9320$, $0.9548$). This is the headline result.
    \item \textbf{all $\approx$ box} --- adding a point to a box adds marginal noise,
          not benefit.
    \item \textbf{Point underperforms text on C$\to$C} ($0.9281 < 0.9329$) --- for thin
          elongated cadaver instruments, the centre-of-mass falls on a background gap.
    \item \textbf{Box on fold2: $+0.093$ Dice} ($0.839\to0.932$) --- largest gain in
          the study.
\end{itemize}

\begin{figure}[H]
    \centering
    \begin{subfigure}[b]{0.32\linewidth}
        \includegraphics[width=\linewidth,height=5.5cm,keepaspectratio]{../../outputs/H4/C_C_fold1_prompt_modebox/predictions/predictions.eps}
        \caption{C$\to$C box --- 0.9476}
    \end{subfigure}
    \begin{subfigure}[b]{0.32\linewidth}
        \includegraphics[width=\linewidth,height=5.5cm,keepaspectratio]{../../outputs/H4/C_C_fold1_prompt_modepoint/predictions/predictions.eps}
        \caption{C$\to$C point --- 0.9281}
    \end{subfigure}
    \begin{subfigure}[b]{0.32\linewidth}
        \includegraphics[width=\linewidth,height=5.5cm,keepaspectratio]{../../outputs/H4/C_C_fold1_prompt_modeall/predictions/predictions.eps}
        \caption{C$\to$C all --- 0.9453}
    \end{subfigure}
    \caption{H4 cadaver: prompt ablation (box / point / all).}
\end{figure}

\begin{figure}[H]
    \centering
    \begin{subfigure}[b]{0.32\linewidth}
        \includegraphics[width=\linewidth,height=5.5cm,keepaspectratio]{../../outputs/H4/L_L_fold1_prompt_modebox/predictions/predictions.eps}
        \caption{L$\to$L fold1 box --- 0.9593}
    \end{subfigure}
    \begin{subfigure}[b]{0.32\linewidth}
        \includegraphics[width=\linewidth,height=5.5cm,keepaspectratio]{../../outputs/H4/L_L_fold1_prompt_modepoint/predictions/predictions.eps}
        \caption{L$\to$L fold1 point --- 0.9132}
    \end{subfigure}
    \begin{subfigure}[b]{0.32\linewidth}
        \includegraphics[width=\linewidth,height=5.5cm,keepaspectratio]{../../outputs/H4/L_L_fold1_prompt_modeall/predictions/predictions.eps}
        \caption{L$\to$L fold1 all --- 0.9580}
    \end{subfigure}
    \caption{H4 live fold1: prompt ablation (box / point / all).}
\end{figure}

\subsection{H5 --- Consolidated Re-evaluation and Failure Analysis}

All 16 H3 and H4 checkpoints were re-tested on held-out test sets. Results match
H3/H4 to 4 decimal places, confirming: (1) deterministic inference,
(2) correct checkpoint selection, (3) no data leakage, (4) end-to-end pipeline
integrity.

Every test run saved a per-image Dice CSV (\texttt{per\_image\_dice.csv}). Analysis
across all 16 runs (23,216 frame evaluations) reveals the failure distribution.

\begin{center}
\begin{tabular}{lrr}
\toprule
\textbf{Criterion} & \textbf{Frames} & \textbf{\% of total} \\
\midrule
Total frames evaluated & 23{,}216 & 100\% \\
Dice $< 0.70$ (weak)   & 1{,}813  & 7.8\% \\
Dice $< 0.50$ (failure)& 1{,}291  & 5.6\% \\
Dice $= 0.00$ (complete miss) & \multicolumn{2}{c}{see worst-frames grid} \\
\bottomrule
\end{tabular}
\end{center}

% GENERATED BY: uv run python scripts/h5_failure_analysis.py
% OUTPUT:       outputs/H5/analysis/dice_distributions.eps
\begin{figure}[H]
    \centering
    \includegraphics[width=\linewidth,height=9cm,keepaspectratio]{../../outputs/H5/analysis/dice_distributions.eps}
    \caption{H5: Left --- per-image Dice box plots across all prompt types (all folds).
             Box prompt consistently shifts the distribution upward with fewer low outliers.
             Right --- Dice histogram for text-prompt runs; 5.6\% of frames fall below 0.5.}
\end{figure}

\textbf{Worst-performing frames} (Dice $= 0.00$, text prompt, cadaver):
The lowest-Dice frames are concentrated in sessions \textbf{S09} and \textbf{S10},
all with Dice $= 0.00$. This indicates \textit{complete segmentation failure} ---
either the instrument is absent from the frame (background-only), or the model
predicts a non-overlapping region. These are not random errors: S09/S10 are specific
cadaver sessions that may have unusual lighting, instrument-absent frames, or
edge cases not well represented in the training split.

% GENERATED BY: uv run python scripts/h5_failure_analysis.py
% OUTPUT:       outputs/H5/analysis/worst_frames_grid.eps
\begin{figure}[H]
    \centering
    \includegraphics[width=\linewidth,height=8cm,keepaspectratio]{../../outputs/H5/analysis/worst_frames_grid.eps}
    \caption{H5: Ten worst-performing frames (text prompt, Dice $< 0.5$).
             All are from cadaver sessions S09/S10. Common patterns: instrument fully
             out of frame, background-only frames, or heavy specular glare.
             These failure modes motivate Week~2 qualitative analysis.}
\end{figure}

\textbf{Implication for Week 2:} The failure frames are already saved --- no new GPU
time is needed. Week~2 will categorise each Dice $< 0.5$ frame into: (a) no
instrument visible, (b) instrument at frame boundary, (c) specular glare/smoke,
or (d) model false positive. This directly explains the fold-to-fold variance
observed in H3 and motivates the prompt design for Phase~2.

% ============================================================
\section{Summary --- Best Results}
% ============================================================

\noindent\textbf{Note on the SOTA reference.}
HFRF-Net (Dice $= 0.9374$) was evaluated on the \textbf{UW-Sinus-Surgery live
subset} --- the same live sessions (L01--L03) used in our L$\to$L experiments.
This makes the L$\to$L comparison \textbf{direct}: same dataset, same split type.
The C$\to$C comparison is indirect (HFRF-Net did not report cadaver results
separately). SAM3+LoRA+box exceeding $0.9374$ on L$\to$L therefore constitutes
a genuine head-to-head improvement over the published task-specific SOTA on this
dataset.

\begin{center}
\begin{tabular}{llcccc}
\toprule
\textbf{Hyp.} & \textbf{Setting} & \textbf{Prompt} & \textbf{Dataset} & \textbf{Dice} & \textbf{$\Delta$ vs ref.} \\
\midrule
H1 & C$\to$C zero-shot      & text & UW-Sinus (cadaver) & 0.6016 & $-0.336$ \\
H2 & L$\to$L zero-shot mean & text & UW-Sinus (live)    & 0.5762 & $-0.361$ \\
H3 & C$\to$C LoRA           & text & UW-Sinus (cadaver) & 0.9329 & $-0.005$ \\
H3 & L$\to$L LoRA mean      & text & UW-Sinus (live)    & 0.877  & $-0.060$ \\
H4 & C$\to$C LoRA           & box  & UW-Sinus (cadaver) & 0.9476 & $+0.010$ \\
\textcolor{sota}{H4} & \textcolor{sota}{L$\to$L fold1} & \textcolor{sota}{box} & \textcolor{sota}{UW-Sinus (live)} & \textcolor{sota}{0.9593} & \textcolor{sota}{$+0.022$} \\
\textcolor{sota}{H4} & \textcolor{sota}{L$\to$L fold2} & \textcolor{sota}{box} & \textcolor{sota}{UW-Sinus (live)} & \textcolor{sota}{0.9320} & \textcolor{sota}{$-0.005$} \\
\textcolor{sota}{H4} & \textcolor{sota}{L$\to$L fold3} & \textcolor{sota}{box} & \textcolor{sota}{UW-Sinus (live)} & \textcolor{sota}{0.9548} & \textcolor{sota}{$+0.017$} \\
\midrule
\multicolumn{4}{l}{HFRF-Net (UW-Sinus live subset, task-specific SOTA)} & 0.9374 & --- \\
\bottomrule
\end{tabular}
\end{center}

\textcolor{sota}{Green}: meets or exceeds HFRF-Net on the same live dataset.
SAM3+LoRA+box beats SOTA on 3 of 4 L$\to$L folds and on C$\to$C (indirect
comparison --- HFRF-Net cadaver not reported). Only L$\to$L fold2, the hardest
session, falls marginally short ($-0.005$).

\subsection*{Direct Comparison on UW-Sinus Live (same data as HFRF-Net)}

% GENERATED BY: uv run python scripts/h5_failure_analysis.py
% INPUT:        outputs/summary.csv
% OUTPUT:       outputs/H5/analysis/direct_comparison.tex
\input{../../outputs/H5/analysis/direct_comparison}

\noindent With text-only prompts (H3), SAM3+LoRA lags HFRF-Net by $-0.060$ on
average --- the gap attributable to the lack of spatial prompt. With a box prompt
(H4), this gap is eliminated: SAM3+LoRA exceeds HFRF-Net in 2 of 3 sessions
($+0.022$, $+0.017$) and nearly matches it in the hardest session ($-0.005$,
fold2). This directly demonstrates that \textbf{spatial prompt quality is the
primary bottleneck}, not the model capacity --- a key motivation for Phase~2
kinematic-derived prompts.

% ============================================================
\section{Discussion}
% ============================================================

\subsection{Why LoRA Works}
Zero-shot SAM3 over-segments because it was pretrained on natural images. LoRA
adapts \texttt{q\_proj}/\texttt{v\_proj} attention to attend to surgical instrument
edges, shifting the model from ``find anything instrument-like'' to tight boundary
delineation --- with only $0.5\%$ trainable parameters.

\subsection{Why Box Dominates}
SAM3 inherits SAM2's training regime (predominantly bounding-box prompts). A box
encodes both location and spatial extent via two corner embeddings. Text gives
semantic identity without spatial constraint; a point gives location but not extent.

\subsection{Why Point Underperforms Text on C$\to$C}
Cadaver instruments are thin and elongated. The centre-of-mass of such shapes
frequently lands on a background gap between instrument parts, providing a
conflicting foreground label at a visually background pixel.

\subsection{Domain Gap}
Cadaver H3 ($0.933$) consistently exceeds live ($0.905$, $0.839$, $0.889$).
Blood, smoke, motion blur, and specular highlights in live surgery degrade
segmentation. High fold-to-fold variance ($\pm0.034$) confirms session-specific
difficulty is the primary factor.

% ============================================================
\section{Research Roadmap}
% ============================================================

\begin{center}
\begin{tabular}{C{1cm} C{4.5cm} C{4cm} C{2.8cm} C{1.2cm}}
\toprule
\textbf{Phase} & \textbf{Question} & \textbf{Data} & \textbf{SAM3?} & \textbf{Weeks} \\
\midrule
1 & Can SAM3+LoRA match SOTA? & UW-Sinus-Surgery-CL (GT prompts) & Done & 1--2 \\
2 & Do kinematic prompts preserve Dice? & EndoVis 2017 (GT masks + video; simulated kinematic noise) & Unchanged & 3--5 \\
3 & Can output feed robot control? & Robot/simulation & Speed opt. & future \\
\bottomrule
\end{tabular}
\end{center}

\subsection{Phase 1 --- Completed}
SAM3+LoRA+box: Dice $0.9476$ (cadaver) and $0.9593$ (live fold1), both exceeding
HFRF-Net. Week~2 adds \textbf{failure mode analysis} using the existing
\texttt{per\_image\_dice.csv} saved per run --- no new GPU time needed.
Worst frames (Dice $<0.5$) categorised by: blood occlusion, instrument at
frame boundary, smoke/haze, specular glare.

\subsection{Phase 2 --- Kinematic-Prompt Robustness on EndoVis 2017 (Weeks 3--5)}

\textbf{Dataset: EndoVis 2017} (MICCAI Robotic Instrument Segmentation Challenge).
10 video clips of abdominal porcine procedures using a da~Vinci robotic system;
300 frames per clip at $1024\times1280$; 7 tool classes with pixel-level GT masks.
Chosen because it provides \textit{both video and GT masks} for controlled evaluation.

\textbf{No public dataset provides video + robot kinematics + GT masks simultaneously.}
JIGSAWS has kinematics but no masks; EndoVis 2017 has masks but no kinematics.
Phase~2 resolves this by \textbf{simulating kinematic error}: a real kinematic prompt
is a noisy estimate of the instrument tip in image space (due to model error,
hand-eye calibration drift, cable friction). We model this directly.

\begin{center}
\begin{tabular}{lp{9.5cm}}
\toprule
\textbf{Step} & \textbf{Detail} \\
\midrule
1. Extract oracle prompt & GT mask $\to$ centroid $(u,v)$ and bounding box ---
                           the \textit{perfect} kinematic estimate \\
2. Inject noise          & $(u,v) \leftarrow (u + \varepsilon,\; v + \varepsilon)$,
                           $\varepsilon \sim \mathcal{N}(0,\,\sigma^2)$;
                           sweep $\sigma \in \{0, 10, 20, 30, 50\}$ px \\
3. Build prompts         & box $= (u{-}r, v{-}r, u{+}r, v{+}r)$; point $= (u,v)$ \\
4. Feed SAM3+LoRA        & Same pipeline as Phase~1; no retraining \\
5. Evaluate              & Dice vs.\ GT mask at each noise level \\
\bottomrule
\end{tabular}
\end{center}

\noindent Each noise level $\sigma$ represents a real-world kinematic scenario:

\begin{center}
\begin{tabular}{cc p{7cm}}
\toprule
$\sigma$ (px) & \textbf{Represents} \\
\midrule
0   & Perfect oracle (= Phase~1 GT prompt) \\
10--20 & Well-calibrated kinematic model \\
30--50 & Poor calibration or kinematic model error \\
\bottomrule
\end{tabular}
\end{center}

\noindent\textbf{Key output:} a Dice-vs-noise robustness curve.
Key metric: $\Delta\text{Dice}(\sigma) = \text{Dice}(\sigma) - \text{Dice}(0)$.
If $|\Delta\text{Dice}| < 0.05$ up to $\sigma = 30$\,px, kinematic prompts are viable
for deployment without requiring sub-pixel calibration accuracy.
This directly answers the deployment question \textit{without needing a robot}.

\subsection{Phase 3 --- Closed-Loop (Future)}
SAM3 output feeds back into the robot controller. Disagreement between kinematic
prediction and SAM3 mask signals tissue contact, occlusion, or calibration drift ---
without force sensors or annotation.

\medskip
\noindent\textit{No prior work combines annotation-free inference + SAM backbone +
kinematic-derived prompts. This gap is the contribution.}

% ============================================================
\section{Repository and Reproducibility}
% ============================================================

\begin{center}
\begin{tabular}{ll}
\toprule
\textbf{Item} & \textbf{Detail} \\
\midrule
Repository  & \texttt{github.com/skhan61/ess-research} \\
Entry point & \texttt{uv run python scripts/run\_all\_experiments.py --plan research\_plan.yaml} \\
Debug       & Add \texttt{--debug} flag \\
Metrics     & \texttt{outputs/\{H\}/\{run\}/metrics.csv} + \texttt{outputs/summary.csv} \\
\bottomrule
\end{tabular}
\end{center}

\textbf{Technical notes:}
\begin{itemize}[noitemsep]
    \item PyTorch 2.6: \texttt{add\_safe\_globals([pathlib.PosixPath])} required for
          checkpoint loading.
    \item SAM3 always requires \texttt{input\_ids} --- no text-free mode.
    \item Spatial prompts normalised manually by $\div$ \texttt{image\_size}.
    \item \texttt{image\_size} must be divisible by 14: use 336, 560, or 1008.
\end{itemize}

\nocite{qin2020uw,hfrfnet2025,ravi2024sam2,rocha2019kinematic,sestini2025safis,yue2024surgicalsam,wu2024pointtracking,liu2024surgsam2}
\bibliographystyle{plain}
\bibliography{../references}

\end{document}
